\documentclass[a4paper,11pt]{article}
\usepackage{exam-style-842-unit}
\examinehead{842 信息化概论 · 第 8 单元练习册}{练习一\quad 电子商务基础与模式}
\begin{document}

\unitcover{842 信息化概论 · 单元练习册}{第 8 单元\quad 电子商务与电子政务}{练习一\quad 电子商务基础与模式（共 100 分）}{约 70 分钟}{电子商务的概念（多方定义、广义与狭义）；电子商务的特点；基于交易主体和交易方式的 B2B、B2C、C2C 模式；B2G、C2B、B2T、O2O 等新兴模式；商业模式与价值链；电子商务平台体系架构（教材未设专节，按大纲补充）。}

\begin{knowmap}
\textbf{本质：}电子商务是依托现代信息技术和网络技术，集金融电子化、管理信息化、商贸网络化于一体，旨在实现物流、资金流与信息流和谐统一的新型贸易方式\\[2pt]
\textbf{广义／狭义：}广义 E-Business（通过电子手段进行的商业事务活动）$\supset$ 狭义 E-Commerce（通过互联网完成的商务交易活动）\\[2pt]
\textbf{四个特点：}业务全球化 · 服务个性化 · 业务集成性 · 机会的均等性\\[2pt]
\textbf{主要模式（按交易主体）：}B2B（企业对企业）· B2C（企业对消费者）· C2C（消费者对消费者）$\Rightarrow$ 新兴 B2G、C2B、BOB、B2T、O2O、ABC、O2P\\[2pt]
\textbf{商业模式与价值链：}商业模式本质是核心业务流程与价值链的有机结合；波特 1985 年《竞争优势》提出价值链＝基本活动＋辅助活动，并扩展为供应链、企业、渠道、买方构成的价值链系统\\[2pt]
\textbf{平台体系架构（按大纲补充）：}基础设施层 $\rightarrow$ 平台服务层（支付、物流、客服、数据等能力组件）$\rightarrow$ 应用业务层（商品、订单、营销、售后）$\rightarrow$ 表示层（Web、App、小程序多端接入），并由安全与运营体系贯穿
\end{knowmap}

\begin{keyterms}
\textbf{电子商务}：依托现代信息技术和网络技术，集金融电子化、管理信息化、商贸网络化于一体，实现物流、资金流与信息流和谐统一的新型贸易方式。\\
\textbf{E-Business}：广义电子商务，指通过电子手段进行的商业事务活动，范围比 E-Commerce 更宽。\\
\textbf{E-Commerce}：狭义电子商务，指通过使用互联网完成的商务交易活动。\\
\textbf{B2B}：企业与企业之间的电子商务，是电子商务应用最多、最受企业重视的形式，代表为阿里巴巴。\\
\textbf{B2C}：企业与消费者之间的电子商务，即企业通过网络向个人消费者销售产品或服务。\\
\textbf{C2C}：消费者与消费者之间的电子商务，平台为买卖双方提供在线交易与竞价环境，代表为 eBay、淘宝。\\
\textbf{B2G}：企业与政府之间的电子商务，涵盖政府采购、税收、商检、政策法规发布等事务。\\
\textbf{C2B}：消费者与企业之间的电子商务，消费者按自身需求定制产品与价格、参与设计生产。\\
\textbf{O2O}：Online to Offline，线上与线下结合的电子商务，让互联网成为线下交易的前台。\\
\textbf{商业模式}：企业为客户创造价值并为自己盈利而开展业务的基本方式，本质是核心业务流程与价值链的有机结合。\\
\textbf{价值链}：原材料经过一系列价值增加活动最终转化为满足客户需要的产品或服务的动态过程。\\
\textbf{价值网}：以网络和信息系统为手段，连接供应商、经销伙伴、分销商与顾客，协同创造价值的网络化模式。\\
\textbf{电子商务平台}：承载商品展示、交易撮合、支付结算、物流配送与数据服务等功能的信息系统与服务体系。\\
\textbf{电子商务平台体系架构}：支撑电商业务的分层结构，通常分为基础设施层、平台服务层、应用业务层与表示层。
\end{keyterms}

\newpage
\begin{contrast}{B2B / B2C / C2C（每行依次为：交易主体·典型代表·主要特点）}
\citem{B2B}{企业与企业}{阿里巴巴}{电子商务应用最多、最受企业重视，可降低价格、提高生产率、降低劳动成本}
\citem{B2C}{企业与消费者}{企业网上商城、网络购物网站}{企业直接向个人销售产品或服务，含虚拟社群、交易聚合、广告网络、O2O 四种经营模式}
\citem{C2C}{消费者与消费者}{eBay、淘宝}{平台提供在线交易与竞价环境，协助汇集市场资讯、建立信用评价制度，平台不负责物流}
\end{contrast}

\begin{contrast}{广义电子商务 / 狭义电子商务（每行依次为：英文缩写·覆盖范围·典型业务）}
\citem{广义电子商务}{E-Business}{通过电子手段进行的商业事务活动，贯穿企业行为全过程}{企业与企业的联系网络、企业内外部协同、供应链与客户关系管理}
\citem{狭义电子商务}{E-Commerce}{通过使用互联网完成的商务交易活动}{网上购物、网上支付、在线订票等具体交易}
\end{contrast}

\begin{contrast}{O2O / C2B / B2G（每行依次为：全称·核心含义·典型场景）}
\citem{O2O}{Online to Offline}{将线下商务机会与互联网结合，让互联网成为线下交易的前台}{线上揽客、到店体验消费、线上结算，推广效果可查、每笔交易可跟踪}
\citem{C2B}{Consumer to Business}{消费者根据自身需求定制产品和价格，或参与产品设计、生产与定价}{个性化定制、反向团购、消费者参与设计}
\citem{B2G}{Business to Government}{企业与政府之间的电子商务，政府既是采购者又是调控与监管者}{政府采购、电子税务、网上商检、政策法规发布}
\end{contrast}

\newpage
\begin{contrast}{商业模式 / 价值链 / 价值网（每行依次为：关注对象·核心内容·相互关系）}
\citem{商业模式}{单个企业如何创造价值并盈利}{产品服务、信息流、资金流、物流、客户收益与效益来源的组织方式}{明确企业在价值链中的战略定位与利润来源}
\citem{价值链}{单个企业的价值活动分解}{波特 1985 年《竞争优势》：基本活动＋辅助活动构成的动态价值创造过程}{是分析企业竞争优势的基本工具}
\citem{价值网}{价值链上多成员的协同网络}{以网络和信息系统连接供应商、分销商与顾客，共享信息、协同创造价值}{价值链柔性化、精细化、适应性增强后形成价值线，进而动态合成价值网}
\end{contrast}

\qtypename{一、单项选择题}{10 题，每题 3 分，共 30 分}
\q 目前较为权威的电子商务概念表述，出自 1997 年在法国巴黎举行的（\quad）
\choicesline{A. 世界电子商务会议}{B. 世界贸易组织会议}{C. 联合国经合组织会议}{D. 全球信息基础设施委员会会议}

\q 广义上讲，电子商务（Electronic Business）是指（\quad）
\choices{A. 通过使用互联网完成的商务交易活动}{B. 通过电子手段进行的商业事务活动}
{C. 仅指网上购物}{D. 仅指电子支付}

\q 狭义电子商务（E-Commerce）是指（\quad）
\choicesline{A. 互联网上的商务交易活动}{B. 所有电子手段的商业事务}{C. 企业内部管理}{D. 政府网上采购}

\q IBM 提出的电子商务定义公式是（\quad）
\choicesline{A. 电子商务＝Web＋IT}{B. 电子商务＝Internet＋EDI}{C. 电子商务＝物流＋资金流}{D. 电子商务＝计算机＋网络}

\q B2B 电子商务的代表模式是（\quad）
\choicesline{A. 阿里巴巴}{B. 淘宝}{C. eBay}{D. 京东}

\q C2C 电子商务的代表是（\quad）
\choicesline{A. 阿里巴巴}{B. eBay、淘宝}{C. 京东}{D. 苏宁易购}

\q O2O 是指（\quad）
\choices{A. 企业与企业之间的电子商务}{B. 线上与线下结合的电子商务}
{C. 消费者与政府之间的电子商务}{D. 企业与政府之间的电子商务}

\q 下列不属于电子商务特点的是（\quad）
\choicesline{A. 业务全球化}{B. 服务个性化}{C. 业务集成性}{D. 交易垄断性}

\q 1985 年在《竞争优势》一书中提出价值链理论的学者是（\quad）
\choicesline{A. 迈克尔·波特}{B. 布鲁克斯}{C. 罗伊斯}{D. 梅棹忠夫}

\q 电子商务平台体系架构中，直接面向用户、负责页面展示与交互的部分通常称为（\quad）
\choicesline{A. 数据层}{B. 表示层（前端）}{C. 基础设施层}{D. 网络层}

\qtypename{二、多项选择题}{4 题，每题 4 分，共 16 分；每题有两个或两个以上正确项}
\q 与传统商务相比，电子商务的特点包括（\quad）
\choicesline{A. 业务全球化}{B. 服务个性化}{C. 业务集成性}{D. 机会的均等性}

\q B2C 电子商务的经营模式包括（\quad）
\choicesline{A. 虚拟社群}{B. 交易聚合}{C. 广告网络}{D. 线上与线下结合（O2O）}

\q 下列属于新兴电子商务模式的有（\quad）
\choicesline{A. B2G}{B. C2B}{C. O2O}{D. B2T}

\q 电子商务平台体系架构通常可以划分为（\quad）
\choices{A. 表示层（Web、App、小程序多端接入）}{B. 应用业务层（商品、订单、支付、物流等）}
{C. 数据层（数据库、数据仓库）}{D. 基础设施层（网络、服务器、存储）}

\qtypename{三、判断题}{6 题，每题 2 分，共 12 分；正确的画“√”，错误的画“×”并写出改正后的表述}
\q 电子商务就是网上购物。\quad（\quad）\quad 改正：\fixblank

\q B2B 是电子商务应用最多、最受企业重视的形式。\quad（\quad）\quad 改正：\fixblank

\q C2C 平台经营者负责商品的物流配送。\quad（\quad）\quad 改正：\fixblank

\q 广义电子商务（E-Business）的范围比狭义电子商务（E-Commerce）更宽。\quad（\quad）\quad 改正：\fixblank

\q 电子商务的机会均等性对中小型企业尤其有利。\quad（\quad）\quad 改正：\fixblank

\q 价值链只是企业内部的活动，与供应商价值链、渠道价值链和买方价值链无关。\quad（\quad）\quad 改正：\fixblank

\qtypename{四、配对题}{6 对，每对 1 分，共 6 分；把 A 栏与 B 栏对应内容配对}
\q 将下列模式、人物或事件与其对应内容配对，把 B 栏字母填入括号。
\par\vspace{0.2em}
\noindent\begin{tabular}{@{}p{0.44\textwidth}@{\hspace{1.2em}}p{0.44\textwidth}@{}}
（1）B2B \hfill（\quad） & A. 阿里巴巴 \\
（2）C2C \hfill（\quad） & B. 淘宝、eBay \\
（3）O2O \hfill（\quad） & C. 线上揽客、线下消费体验 \\
（4）世界电子商务会议 \hfill（\quad） & D. 1997 年巴黎，国际商会 \\
（5）迈克尔·波特 \hfill（\quad） & E. 1985 年《竞争优势》与价值链 \\
（6）IBM \hfill（\quad） & F. 电子商务＝Web＋IT \\
\end{tabular}\par\vspace{0.3em}

\qtypename{五、简答题}{4 题，每题 6 分，共 24 分}
\q 什么是电子商务？比较广义电子商务与狭义电子商务，并给出一个较为完整的综合定义。\anslines{2.4cm}

\q 简述电子商务的主要特点。\anslines{2.4cm}

\q 简述按交易主体和交易方式划分的 B2B、B2C、C2C 三种电子商务模式及其典型代表。\anslines{2.4cm}

\q 简述电子商务平台体系架构的层次划分及各层作用。\anslines{2.4cm}

\newpage
\qtypename{六、综合题}{1 题，共 12 分}
\q 某县依托电商平台销售猕猴桃：平台上有企业自营旗舰店，有合作社在平台上开店，也有消费者之间的二手交易；同时开展“线上下单、到果园体验采摘”活动，并与县政府合作参与消费帮扶采购。请回答：
（1）分别判断企业自营旗舰店、合作社开店、消费者二手交易、线上引流线下采摘、政府消费帮扶采购各属于哪种电子商务模式；
（2）用价值链与价值网说明电商平台如何为农户、合作社和消费者创造价值；
（3）为该县设计一个电子商务平台体系架构的层次框架，并说明各层的作用。\anslines{5cm}

\end{document}
